\section{Background and Related Work}
\label{sec:background}

\subsection{Photographic Film as an Imaging System}

A color negative film is a stack of light-sensitive layers coated on a
transparent base. In simplified form, from the incident-light side inward:
an ultraviolet-absorbing overcoat; a blue-sensitive layer forming yellow dye; a
yellow filter layer that prevents residual blue light from reaching the layers
below; a green-sensitive layer forming magenta dye; a red-sensitive layer
forming cyan dye; the acetate or polyester base; and an antihalation backing
that absorbs light which passes entirely through the emulsion.

Each sensitized layer contains silver halide microcrystals suspended in
gelatin. When a crystal absorbs sufficient photons, a small cluster of metallic
silver atoms forms on its surface --- the \emph{latent image}. Development
chemically reduces the entire crystal to metallic silver, an amplification of
order $10^8$--$10^9$. In color film, the oxidized developer then couples with
dye-forming compounds in the layer, producing dye clouds; the silver is
subsequently bleached and fixed away, leaving only dye.

Four consequences of this architecture drive the entire \EM{} model:

\begin{enumerate}
\item \textbf{The response is logarithmic and sigmoidal.} Because
crystals have a distribution of sizes and sensitivities, the fraction developed
as a function of exposure follows a smooth cumulative distribution, producing
the characteristic $D$ versus $\logE$ curve with its toe, straight line, and
shoulder. This is modelled in Section~\ref{sec:negative}.
\item \textbf{The image is granular by construction.} The final density is
the sum of discrete dye clouds at effectively random positions. The signal is
not merely noisy; its noise is structurally tied to how much of the image was
formed. This is modelled in Section~\ref{sec:grain}.
\item \textbf{Layers interact.} Development inhibitors released by one layer
diffuse laterally and into adjacent layers, suppressing development there. This
is the DIR (Development Inhibitor Releasing) coupler mechanism, and it is
responsible for both the edge-enhancement that gives negative film its apparent
acutance and the cross-channel behavior that keeps saturation controlled at
high density. This is modelled in Section~\ref{sec:interlayer}.
\item \textbf{Light scatters inside the stack.} Light not absorbed on first
pass reflects from layer interfaces and from the base, re-exposing the emulsion
at a displaced position. Because the red-sensitive layer sits deepest and red
light penetrates furthest, the effect is strongly red-biased. This is halation,
modelled in Section~\ref{sec:halation}.
\end{enumerate}

The negative is then optically printed: white light passes through the negative
and exposes a \emph{print stock}, which has its own, much steeper,
characteristic curve. The print stock is where most of the final contrast and
color character is created. Reproducing this second transfer is what separates
a ``film look'' from a ``scan of a negative'' look, and its omission is the
most common failure of consumer film simulations.

\subsection{Prior Art in Professional Tools}

Desktop film emulation divides into three technical families.

\subsubsection{Lookup-table pipelines}
The dominant approach in color grading practice applies a 3D lookup table (3D
LUT), typically $33^3$ or $65^3$, generated offline. Print film emulation LUTs
are the canonical example: a single LUT captures the composite negative-plus-print
transfer for a fixed set of printer lights. This is fast and reproducible, and
it is genuinely correct for the pointwise part of the transform. Its limitation
is definitional: a LUT cannot carry spatial or stochastic behavior, so halation,
grain, and interlayer effects must be added by separate operators, and the LUT
must be regenerated whenever a parameter inside the baked chain changes. \EM{}
uses LUTs as an \emph{optimization} for the pointwise sub-chain
(Section~\ref{sec:lutbake}) but never as the model.

\subsubsection{Parametric density models}
Tools in the Dehancer and FilmConvert family expose parameters that plausibly
correspond to film behavior and evaluate a parametric transfer function per
pixel, augmented by spatial passes for halation and bloom and a procedural
grain generator. This is the family \EM{} belongs to. The distinguishing
questions within the family are how faithfully the density model is fit to
sensitometry, whether the grain model is statistically grounded or merely
textural, and whether the print stage is a separate physical transfer or a
cosmetic contrast curve.

\subsubsection{Spectral and layer-resolved simulation}
The most rigorous approach models the spectral sensitivity of each layer, the
spectral transmittance of each dye, and the spectral power distribution of the
printer illuminant, integrating over wavelength rather than over three
channels. This handles metamerism and illuminant changes correctly and predicts
crosstalk from first principles rather than by fitted matrix. It is
computationally expensive --- typically 31 or more wavelength samples --- and
requires spectral data that is not publicly available for most stocks. \EM{}
adopts a \emph{reduced spectral} strategy (Section~\ref{sec:colorspace}): the
crosstalk matrices and the printer-light response are derived from a spectral
model offline, at profile authoring time, and baked into three-channel
operators for runtime. This captures most of the spectral benefit at
three-channel cost, and it leaves the door open to a full spectral path in a
later release.

\subsection{Prior Art in Mobile Applications}

Table~\ref{tab:priorart} summarizes the capability boundary. The column
\emph{Scene-referred} is the decisive one: applications that begin from a
processed, display-referred image have already lost the highlight information
that the film shoulder is supposed to compress, and they cannot recover it.

\begin{table}[htbp]
\caption{Capability Comparison, Mobile Film Simulation}
\label{tab:priorart}
\begin{center}
\renewcommand{\arraystretch}{1.15}
\begin{tabular}{@{}lccccc@{}}
\toprule
\textbf{Class} & \textbf{Scene} & \textbf{Print} & \textbf{Proc.} & \textbf{Phys.} & \textbf{Non-} \\
 & \textbf{ref.} & \textbf{stage} & \textbf{grain} & \textbf{halation} & \textbf{destr.} \\
\midrule
Retro filter apps      & --- & --- & --- & --- & --- \\
LUT-based camera apps  & --- & $\circ$ & --- & --- & $\circ$ \\
Pro mobile RAW editors & \checkmark & --- & $\circ$ & --- & \checkmark \\
Desktop film emulation & \checkmark & \checkmark & \checkmark & \checkmark & \checkmark \\
\midrule
\textbf{\EM{}}         & \checkmark & \checkmark & \checkmark & \checkmark & \checkmark \\
\bottomrule
\multicolumn{6}{@{}l@{}}{\footnotesize \checkmark\ full; $\circ$\ partial or cosmetic; ---\ absent.}
\end{tabular}
\end{center}
\end{table}

The observation motivating \EM{} is that the rightmost row of
Table~\ref{tab:priorart} has, until recently, been unreachable on mobile for
performance reasons. That constraint has lapsed. An A17-class or later mobile
GPU sustains several teraflops of half-precision throughput with unified memory
and tile-local storage; the pipeline specified in Section~\ref{sec:metal}
requires approximately 40--60 arithmetic operations and 3 texture fetches per
pixel for the pointwise chain, plus four separable spatial passes. At
$4032\times3024$ this is a workload measured in tens of milliseconds, not
seconds. The gap in the market is therefore an engineering-effort gap, not a
silicon gap --- which is precisely the kind of gap a focused team can close.

\subsection{Position Relative to Apple's Imaging Stack}

Apple provides three relevant subsystems, and \EM{} deliberately uses each at a
different level of trust:

\begin{itemize}
\item \textbf{AVFoundation / ProRAW capture} is used fully and without
modification. Apple's Bayer demosaic, sensor noise reduction, and lens shading
correction are excellent and reimplementing them would be value-destructive.
\item \textbf{Core Image RAW processing} (\texttt{CIRAWFilter}) is used as the
demosaic and linearization front end, configured to produce a
\emph{scene-referred linear} output with tone mapping, local contrast, and
sharpening disabled. This is the critical configuration decision of the entire
application and is specified in detail in Section~\ref{sec:raw-config}.
\item \textbf{Core Image's filter graph} is \emph{not} used for the development
pipeline. Its concatenation and tiling behavior are opaque, its precision
guarantees are insufficient for a density-domain pipeline, and its
stochastic-noise support cannot express the grain model. The pipeline is
implemented directly in Metal with an application-owned render graph
(Section~\ref{sec:rendergraph}).
\end{itemize}
